\chapter{Praktická část}

\section{Popis dat}

Poskytovatelem dat je Datové centrum Ústeckého kraje (DCUK), které spravuje a zpřístupňuje data z různých systémů veřejné správy.
Použitý dataset má velikost přibližně 650 MB a obsahuje zhruba 600 tisíc záznamů. Data jsou ve strukturovaném formátu CSV.
Jedná se o záznamy ze systému pro sledování pohybu vozidel správy a údržby silnic. 
Každý záznam obsahuje zejména identifikátor vozidla, typ aktivity a geografickou polohu. 
Tyto informace umožňují analýzu pohybu vozidel v čase a prostoru, například pro vyhodnocení efektivity údržby silniční sítě nebo sledování tras jednotlivých vozidel.

\section{Předběžná srovnávací analýza systémů}

% proč to vzniklo, kdy to zvzniklo, lehká historie
\subsection{ClickHouse}

ClickHouse je sloupcově orientovaná databáze určená pro analytické dotazy v reálném čase (OLAP), původně vyvinutá společností Yandex jako interní nástroj pro zpracování velkých objemů dat. 
Později byla uvolněna jako open-source projekt a postupně se stala jednou z populárních databází pro analytické zpracování dat. \cite{clickhouseintro}
Hlavním důvodem vzniku ClickHouse byla potřeba efektivního zpracování a analýzy rozsáhlých logových a telemetrických dat v reálném čase, kde tradiční relační databáze nedokázaly splnit požadavky na výkon.

ClickHouse je dostupný jak ve formě open-source řešení pro self-hosting, tak i jako plně spravovaná cloudová služba. Open-source varianta umožňuje nasazení na vlastních serverech nebo v libovolném cloudovém prostředí.

\subsection{MariaDB ColumnStore}

MariaDB ColumnStore je sloupcové rozšíření relační databáze MariaDB zaměřené na analytické zpracování velkých objemů dat. 

ColumnStore umožňuje distribuované zpracování dat a paralelizaci dotazů, což výrazně zlepšuje výkon při práci s rozsáhlými datasetty. 
Je vhodný zejména pro analytické reporty a agregace nad velkými tabulkami. \cite{mariadbcolumnstore}

\subsection{InfluxDB}

InfluxDB je časově orientovaná databáze navržená pro ukládání a analýzu dat závislých na čase, například metrik, senzorických dat nebo logů.

Vznikla v roce 2013 společností InfluxData jako odpověď na rostoucí potřebu efektivního ukládání a dotazování časových řad. 
Je optimalizována pro vysokou rychlost zápisu a práci s časovými intervaly, což ji činí vhodnou pro IoT, monitoring a observabilitu systémů. \cite{influxdbdocs}

\subsection{IoTDB}

IoTDB je časově orientovaná databáze vyvinutá Apache Software Foundation zaměřená na IoT scénáře a edge computing.

Byla navržena pro efektivní ukládání a zpracování velkého množství senzorických dat s vysokou frekvencí zápisu. 
Podporuje hierarchickou strukturu dat a je optimalizována pro nízké nároky na úložiště i výpočetní výkon. \cite{iotdbdocs}

\subsection{MariaDB}

MariaDB je relační databázový systém vzniklý jako fork MySQL po jeho převzetí společností Oracle v roce 2009. Cílem projektu bylo zachovat open-source charakter a komunitní vývoj.

MariaDB je široce používaná pro transakční systémy, webové aplikace a běžné ukládání strukturovaných dat. 
Nabízí vysokou kompatibilitu s MySQL a stabilní výkon pro klasické databázové operace. \cite{mariadbdocs}

\subsection{MongoDB}

MongoDB je dokumentově orientovaná NoSQL databáze, která ukládá data ve formátu BSON (binární JSON). 
Byla navržena pro flexibilní práci se schématy a snadnou škálovatelnost.

Je vhodná pro aplikace, kde se struktura dat často mění, například webové služby, katalogy nebo IoT aplikace. \cite{mongodbdocs}

\subsection{QuestDB}

QuestDB je vysoce výkonná open-source time-series databáze zaměřená na nízkou latenci a vysokou propustnost při zpracování dat v reálném čase.

Byla navržena s důrazem na jednoduchost, SQL kompatibilitu a extrémně rychlé dotazy nad časovými řadami. 
Využívá se zejména v oblasti finančních dat, monitoringu a IoT. \cite{questdbdocs}

\subsection{TimescaleDB}

TimescaleDB je rozšíření relační databáze PostgreSQL zaměřené na práci s časovými řadami. Kombinuje klasický SQL přístup s optimalizacemi pro time-series data.

Vznikla v roce 2017 jako open-source projekt s cílem propojit výhody relačního modelu a specializovaných time-series databází. 
Je vhodná pro monitoring, IoT a analytické systémy, kde je potřeba jak flexibilita SQL, tak efektivní práce s časem. \cite{timescaledbdocs}
 
\myworries{Dopsat informace k jednotlivým databázím}

\section{Návrh testovacího scénáře}

Testovací prostředí je postaveno na virtualizační technologii LXD, která umožňuje vytváření izolovaných kontejnerů s 
nízkou režijní náročností. Každý testovaný databázový systém běží ve svém samostatném LXD kontejneru, čímž je 
zajištěna oddělenost prostředí a reprodukovatelnost výsledků. Všechny kontejnery mají přidělené stejné hardwarové prostředky, aby 
byly výsledky jednotlivých testů mezi databázemi porovnatelné.

Proces vytváření kontejnerů, instalace databázových systémů a spuštění testů je plně automatizován pomocí 
nástroje Ansible. Tím je zajištěna reprodukovatelnost testovacího prostředí napříč všemi testy a minimalizace lidského zásahu.

Struktura Ansible projektu je navržena modulárně:

\begin{itemize}
    \item \textbf{Vytvoření kontejneru} \\
    Pomocí Ansible modulu \texttt{community.general.lxd\_container} je vytvořen nový LXD kontejner se specifikovanými parametry.
    
    \item \textbf{Instalace databázového systému} \\
    Po vytvoření kontejneru následuje instalace zvolného databázového systému podle odpovídající Ansible role. 
    Každý systém má vlastní roli, která zajišťuje instalaci, konfiguraci a spuštění služby.
    
    \item \textbf{Příprava na testování} \\
    Po instalaci databáze je kontejner nakonfigurován pro testování. 
    Nahrájí se testovací skripty, dataset a konfigurační soubory.
\end{itemize} 

% připravit archi schéma

Cílem testovacího scénáře je porovnat vybrané databázové systémy z hlediska výkonu ve vybraných testech. Testy budou zaměřeny na následující oblasti:
\begin{itemize}
    \item Rychlost jednorázového zápisu dat 
    \item Vrácení jednoho specifického záznamu
    \item Vrácení všech záznamů za určitý časový interval
    \item Vrácení všech záznamů splňujících určitou podmínku
    \item Agregační dotazy
    \item Rychlost zápisu dat v reálném čase
\end{itemize}

Během všech testů se budou sbírat metriky jako doba trvání dotazu, využití CPU, využití paměti a I/O operace disku.
Všechny tyto metriky budou následně zapsány přes Node Exportery v každém virtuálním stroji.
Ty budou následně dány k dispozici službě Prometheus, která bude zajišťovat sběr a uchovávání těchto metrik pro následnou analýzu a vizualizaci Grafanou.

Veškeré testy se budou provádět na stejném hardwarovém vybavení a za stejných podmínek, na virtuálních strojích s předem definovanými parametry (CPU, RAM, diskový prostor).


\section{Konfigurace testovacího prostředí a databázových systémů}

\subsection{Virtualizace}
Cílem bylo vytvořit testovací prostředí, které bude co nejvíce izolované a zároveň umožní snadnou správu a automatizaci.
Z těchto důvodů jsem si zvolil virtualizační technologii LXD, která umožňuje vytváření a správu kontejnerů a virtuálních strojů.

Zvolil jsem si virtuální stroje pro jejich větší izolaci a lepší kontrolu nad prostředím, ve kterém budou testy probíhat.
To zajistí, že výsledky testů budou co nejvíce objektivní a porovnatelné mezi jednotlivými databázovými systémy.

Další výhodou tohoto přístupu je možnost snadno resetovat testovací prostředí mezi jednotlivými testy, což zajišťuje konzistenci a reprodukovatelnost výsledků.
\subsection{Architektura}
Každý testovaný databázový systém běží ve svém samostatném LXD kontejneru, který je nakonfigurován s přidělenými hardwarovými prostředky (CPU, RAM, diskový prostor).

Při vytvoření kontejneru se do něj nainstaluje zvolený databázový systém a Node Exporter pro sběr metrik.

Na vyhrazeném virtuálním zdroji běží Grafana a Prometheus, které zajišťují sběr a vizualizaci metrik z jednotlivých kontejnerů.

\subsection{Ansible}
\begin{lstlisting}[caption={Struktura adresářů pro Ansible}, label={lst:ansiblestructure}, captionpos=b]
    root/
    |
    |-- ansible.cfg
    |-- setup.sh
    |-- run_test.sh
    |-- run_live_insert.sh
    |
    |-- setup/
    |   |-- setup.yml
    |   |-- setup_monitoring.yml
    |   \-- reload_prometheus.yml
    |
    |-- vars/
    |   |-- resource_profiles.yml
    |   \-- test_types.yml
    |
    |-- roles/
    |   |-- common-create/
    |   |-- common-start/
    |   |-- common-post-processing/
    |   |-- profile-loader/
    |   |-- node-exporter/
    |   |-- monitoring/
    |   |   |-- tasks/
    |   |   |-- templates/
    |   |   \-- files/
    |   |
    |   |-- clickhousedb/
    |   |   |-- tasks/
    |   |   |   |-- install.yml
    |   |   |   \-- setup_database.yml
    |   |   \-- defaults/
    |   |       \-- main.yml
    |   |
    |   |-- ...
    |   \-- columnstoredb/
    |
    \-- tests/
        |-- clickhousedb/
        |   |-- insert.yml
        |   |-- aggregate.yml
        |   |-- condition.yml
        |   |-- single_point.yml
        |   \-- live_insert.yml
        |
        |-- ...
        |-- columnstoredb/
        \--  generate_live_data.py
\end{lstlisting}
Projekt se skládá z několika klíčových adresářů a souborů:
\begin{itemize}
    \item \textbf{inventory/} \\
    Inventář obsahuje definice hostů, databázových systémů, které mají přednastavenou IP adresu.
    \item \textbf{roles/} \\
    Role v projektu bych rozdělil do 3 kategorií:
    \begin{itemize}
        \item \textbf{common} \\
        Obsahuje role pro vytváření kontejnerů, jejich spuštění a post-processing výsledků.
        \item \textbf{monitoring} \\
        Role pro instalaci a konfiguraci Node Exporterů, Prometheus a Grafany.
        \item \textbf{databases} \\
        Každý testovaný databázový systém má svou vlastní roli, která obsahuje úkoly pro instalaci, konfiguraci a spuštění služby.
        Dále obsahuje i datový soubor, který následně testovací skripty použijí pro naplnění databáze testovacím datasetem.
    \end{itemize}
    \item \textbf{tests/} \\
    Adresář obsahuje definice testů pro jednotlivé databázové systémy.
    \item \textbf{setup/} \\
    Adresář obsahuje ansible playbooku, které slouží jako vstupní body pro vytvoření testovacího prostředí a nastavení monitorovacího systému.
\end{itemize}

\subsection{Konfigurace databází}
% pospat jak se konfiguruje LXD
% speciální konfigurace DB

\section{Pilotní nasazení systémů}

Pilotní testovaní bylo provedeno na mém osobním notebooku, který má následující hardwarové parametry:
\begin{itemize}
    \item AMD Ryzen™ 5 5600H (max. zvýšený takt 4,2 GHz, 16 MB mezipaměti L3, 6 jader, 12 vláken)
    \item 16GB paměť DDR4-3200 MHz RAM (2x 8 GB)
    \item NVIDIA® GeForce® GTX 1650 (4 GB vyhrazené paměti GDDR6)
    \item Disková jednotka SSD 512 GB PCIe® NVMe™ M.2
    \item Integrovaná síťová karta 10/100/1000 GbE LAN
\end{itemize}
Na notebooku je nainstalován operační systém Ubuntu 22.04 LTS v minimální (serverové) konfiguraci bez grafického uživatelského rozhraní. 
Systém je ovládán výhradně prostřednictvím příkazové řádky (CLI).
Minimální konfiguraci jsem zvolil z důvodu optimalizace výkonu a minimalizace spotřeby systémových prostředků, které by mohly ovlivnit výsledky testů.

% popsat proxmox vpn
% popsat vpn
% popsat vpn hardware

\section{Implementace testů}
Pro testování databázových systémů jsem následující sadu testů:
\begin{itemize}
    \item \textbf{Insert test} \\
    Test zaměřený na měření rychlosti zápisu dat do databáze.
    \item \textbf{Single point test} \\
    Test zaměřený na měření rychlosti vrácení jednoho specifického záznamu z databáze.
    \item \textbf{Condition test} \\
    Test zaměřený na měření rychlosti vrácení všech záznamů splňujících určitou podmínku.
    Jako podmínku jsem si zvolil vrácení všech záznamů, které mají polohu v určitém geografickém rozsahu za určité období.
    \item \textbf{Aggregate test} \\
    Test zaměřený na měření rychlosti provedení agregačního dotazu nad databází.
    V agregačním testu jsem zvolil dotaz, který vrací průměrnou rychlost, maximální rychlost a počet záznamů pro každé vozidlo za určité období.
\end{itemize}
Všechny tyto testy se spouští pomocí skriptu \texttt{run\_test.sh}, který přijímá parametry pro výběr databázového systému a typu testu.
\begin{lstlisting}[caption={Příklad spuštění live insert testu}, label={lst:liveinsert}, captionpos=b]
    ./run_test.sh -d clickhousedb -t insert
\end{lstlisting}
\begin{itemize}
     \item \textbf{Live insert test} \\
    Test zaměřený na měření rychlosti zápisu dat v reálném čase, kdy se do databáze kontinuálně vkládají nové záznamy.
    Má své speciální parametry:
    \begin{itemize}
        \item \textbf{workers} \\
        Počet paralelních zdrojů, které vkládají data do databáze.
        \item \textbf{duration} \\
        Celková doba v sekundách, po kterou bude test probíhat.
        \item \textbf{batch\_size} \\
        Počet záznamů v jedné dávce, která bude vložena do databáze najednou.
    \end{itemize}
    Test se zavolá následujícím způsobem:
    \begin{lstlisting}[caption={Příklad spuštění live insert testu}, label={lst:liveinsert}, captionpos=b]
    ./run_live_insert.sh -d clickhousedb -w 4 -t 300 -b 1000
    \end{lstlisting}
    Příkaz nastavý test pro databázový systém ClickHouse přes parametr \texttt{-d}, počet zdrojů na 4 přes parametr \texttt{-w}, dobu trvání testu na 300 sekund (5 minut) přes parametr \texttt{-t} a velikost dávky na 1000 záznamů přes parametr \texttt{-b}.
\end{itemize}

Ze všech testů se sbírají metriky dvojím způsobem:
\begin{itemize}
    \item \textbf{NDJSON výsledky} \\
    Po dokončení každého testu se výsledky ukládají do souboru ve formátu NDJSON (Newline Delimited JSON), který obsahuje informace o době trvání dotazu, počtu záznamů vrácených dotazem...
    \item \textbf{Prometheus metriky} \\
    Během testů se v každém kontejneru sbírají metriky pomocí Node Exporteru, které jsou následně shromažďovány službou Prometheus. 
    Tyto metriky zahrnují využití CPU, využití paměti, I/O operace disku a síťovou aktivitu.
\end{itemize}


